\documentclass[11pt]{article}

% ── packages ──────────────────────────────────────────────────────────────────
\usepackage[margin=1.2in]{geometry}
\usepackage{amsmath,amssymb,amsthm}
\usepackage{mathtools}
\usepackage{bm}
\usepackage{hyperref}
\usepackage{enumitem}
\usepackage{microtype}
\usepackage{xcolor}
\usepackage{booktabs}
\usepackage{array}
\usepackage{parskip}
\usepackage{natbib}

% load cleveref AFTER hyperref
\usepackage[capitalise,noabbrev]{cleveref}

% ── hyperref setup ────────────────────────────────────────────────────────────
\hypersetup{
  colorlinks=true,
  linkcolor=blue!60!black,
  citecolor=blue!60!black,
  urlcolor=blue!60!black,
  pdftitle={Near-Orthogonality of Layer Jacobians in Deep Residual Networks},
  pdfauthor={Amir Hossein Rasti}
}

% ── theorem environments ──────────────────────────────────────────────────────
\theoremstyle{plain}
\newtheorem{theorem}{Theorem}[section]
\newtheorem{lemma}[theorem]{Lemma}
\newtheorem{corollary}[theorem]{Corollary}
\newtheorem{proposition}[theorem]{Proposition}

\theoremstyle{definition}
\newtheorem{assumption}{Assumption}
\newtheorem{definition}[theorem]{Definition}

\theoremstyle{remark}
\newtheorem{remark}[theorem]{Remark}

% tell cleveref how to name the assumption counter
\crefname{assumption}{Assumption}{Assumptions}
\Crefname{assumption}{Assumption}{Assumptions}

% ── notation macros ───────────────────────────────────────────────────────────
\newcommand{\R}{\mathbb{R}}
\newcommand{\E}{\mathbb{E}}
\newcommand{\Prob}{\mathbb{P}}
\newcommand{\N}{\mathcal{N}}
\newcommand{\smin}{\sigma_{\min}}
\newcommand{\smax}{\sigma_{\max}}
\newcommand{\norm}[1]{\left\|#1\right\|}
\newcommand{\abs}[1]{\left|#1\right|}
\newcommand{\inner}[2]{\left\langle #1,\, #2 \right\rangle}
\newcommand{\tr}{\operatorname{tr}}
\newcommand{\rank}{\operatorname{rank}}
\newcommand{\reff}{r_{\mathrm{eff}}}
\newcommand{\dI}{d_{\mathcal{I}}}
\newcommand{\Cov}{\operatorname{Cov}}

% ── title block ───────────────────────────────────────────────────────────────
\title{%
  \textbf{Near-Orthogonality of Layer Jacobians\\[0.2em]
  in Deep Residual Networks}\\[0.6em]
  \large A Random-Matrix Proof with Explicit Constants
}
\author{Amir Hossein Rasti}
\date{June 2026}

% ─────────────────────────────────────────────────────────────────────────────
\begin{document}
\maketitle
\vspace{-1em}
\noindent\rule{\linewidth}{0.4pt}
\vspace{0.5em}
\thispagestyle{empty}

% ── abstract ──────────────────────────────────────────────────────────────────
\begin{abstract}
We prove that in an $L$-layer residual network satisfying a depth-scaled
spectral-norm bound on residual block Jacobians, the per-layer contributions
to the full parameter Jacobian are asymptotically \emph{near-orthogonal} in
the Frobenius inner-product sense.  Concretely, for distinct layers
$l \neq l'$,
\[
  \bigl|\inner{J_{\theta}^{(l)}}{J_{\theta}^{(l')}}_F\bigr|
  \;\leq\;
  e^{2c}\,C_{\mathrm{HW}}\,\sqrt{\dI\cdot W}
  \quad\text{with probability}\;\geq 1-2e^{-\dI},
\]
while each diagonal term satisfies
$\norm{J_{\theta}^{(l)}}_F^2 = \Theta\!\left(e^{2c}\,\dI W/L\right)$.
The cross-to-diagonal ratio is therefore $O\!\left(L/\sqrt{\dI W}\right)$,
which vanishes as $W\to\infty$ with $L$ and $\dI$ fixed.
As a corollary, the effective rank of the Gauss--Newton matrix
$G = J_\theta^\top J_\theta$ scales as
$\Omega(L\cdot W_{\min}/\dI)$, growing multiplicatively with both depth
and width.  This closes the rank-summation gap present in prior treatments
of the loss landscape of overparameterised residual networks and provides
the first explicit constant for the cross-term bound.
The proof combines the Hanson--Wright inequality \citep{rudelson2013}
with a product-integral representation of the residual pushforward operator.
\end{abstract}

% ─────────────────────────────────────────────────────────────────────────────
\section{Introduction}
\label{sec:intro}

A central empirical puzzle of deep learning is that gradient-based
optimisers rarely become trapped at suboptimal local minima, even for
highly non-convex loss functions.  For overparameterised networks, several
lines of work establish that the \emph{existence} of bad local minima is
structurally improbable, typically by showing that the Gauss--Newton matrix
$G = J_\theta^\top J_\theta$ carries a high-rank positive-semidefinite
component that dominates the residual curvature
\citep{kawaguchi2016,du2019,oymak2020}.

These treatments share a common gap.  When the full parameter Jacobian
$J_\theta$ is assembled by stacking per-layer contributions
$\{J_\theta^{(l)}\}_{l=1}^{L}$, the effective rank of $G$ is bounded not
by the \emph{sum} of per-layer ranks but by the rank of their union.
The difference is decisive: if the per-layer Jacobians are highly
correlated---sharing optimisation directions---depth provides no advantage
beyond parameter count.  If they are near-orthogonal, ranks sum, and depth
multiplies rather than merely adds to optimisation capacity.

This paper proves the near-orthogonality claim rigorously and provides
explicit constants.  The result is stated for deep residual networks
satisfying a standard depth-scaling assumption on residual block Jacobians
(Assumption~\ref{asm:spectral}), which is implied by gradient clipping or
weight normalisation---both ubiquitous in practice.

\paragraph{Main result (informal).}
For any two layers $l\neq l'$ in an $L$-layer overparameterised residual
network, the Frobenius inner product of their parameter Jacobians satisfies
\[
  \bigl|\inner{J_\theta^{(l)}}{J_\theta^{(l')}}_F\bigr|
  \;=\;
  O\!\Bigl(\sqrt{\dI\cdot W}\,/\,L\Bigr)
\]
with high probability over initialisation, where $\dI$ is the intrinsic
dimension of the data manifold and $W$ is the layer width.  Each diagonal
term is $\Theta(\dI W/L)$, so the cross-term is $o$ of the diagonal by a
factor of $1/\sqrt{\dI W}$.  In the overparameterised regime $W\gg \dI$,
this ratio is small, confirming near-orthogonality.

\paragraph{Implications.}
Near-orthogonality implies that the effective rank of $G$ scales as
$\Omega(L\cdot W_{\min}/\dI)$.  This is strictly stronger than previous
depth-independent bounds: depth and width contribute \emph{multiplicatively}
to optimisation capacity, consistent with empirical observations that deep
residual networks outperform shallow wide ones beyond what parameter count
alone predicts.

\paragraph{Organisation.}
Section~\ref{sec:setup} states assumptions and notation.
Section~\ref{sec:main} gives the complete proof.
Section~\ref{sec:corollary} derives the Gauss--Newton rank corollary.
Section~\ref{sec:scope} discusses scope, limitations, and falsifiable
predictions.

% ─────────────────────────────────────────────────────────────────────────────
\section{Setup and Assumptions}
\label{sec:setup}

\subsection{Architecture}

We consider an $L$-layer residual network with uniform hidden width $W$.
Layer $l$ computes
\[
  x_l \;=\; x_{l-1} + F_l(x_{l-1};\,\theta_l),
  \qquad l=1,\ldots,L,
\]
where $F_l:\R^W\to\R^W$ is the residual block with parameters
$\theta_l\in\R^p$, $p=W^2$ (one weight matrix per layer for clarity; the
result extends immediately to multi-layer blocks).  The output is
$\hat{y} = h(x_L)$ for a fixed readout $h$.  The loss is
$\mathcal{L}(\theta) = \tfrac{1}{n}\sum_{i=1}^n \ell(\hat{y}(x_i),y_i)$
over $n$ data points.

\subsection{Assumptions}

\begin{assumption}[Depth-scaled spectral norm]
\label{asm:spectral}
There exists a universal constant $c>0$ such that for all $l$ and all
$\theta$ along the training trajectory,
\[
  \bigl\|J_{F_l}(\theta)\bigr\|_2 \;\leq\; \frac{c}{L},
\]
where $J_{F_l} = \partial F_l/\partial x_{l-1}$ is the activation Jacobian
of the residual block.  This is enforced in practice by gradient clipping
\citep{zhang2020} or spectral normalisation \citep{miyato2018} and is
consistent with the mean-field initialisation regime.
\end{assumption}

\begin{assumption}[Data excitation]
\label{asm:excitation}
The empirical covariance of representations $x_{l-1}$ satisfies
$\lambda_{\min}\!\left(\Cov(x_{l-1})\right)\geq\gamma>0$ uniformly in $l$.
The data thus activates all $\dI$ intrinsic dimensions at every layer.
\end{assumption}

\begin{assumption}[Sub-Gaussian initialisation]
\label{asm:init}
At initialisation, each weight matrix $W_l$ has i.i.d.\ entries distributed
as $\mathcal{N}(0,2/W)$ (He initialisation).  More generally, it suffices
that entries be mean-zero sub-Gaussian with parameter $\sigma^2=O(1/W)$.
\end{assumption}

\subsection{Key Objects}

\paragraph{Per-layer parameter Jacobian.}
The contribution of layer $l$'s parameters to the output Jacobian is
\[
  J_\theta^{(l)}
  \;=\;
  \Phi_{L\leftarrow l}\cdot Z_l
  \;\in\;
  \R^{n_{\mathrm{out}}\times p},
\]
where $Z_l = \partial x_l/\partial\theta_l\in\R^{W\times p}$ is the
\emph{local parameter Jacobian} and
\[
  \Phi_{L\leftarrow l}
  \;=\;
  \prod_{k=l}^{L-1}\bigl(I + J_{F_k}\bigr)
\]
is the \emph{pushforward operator} mapping layer-$l$ feature space to the
output.

\paragraph{Full Jacobian and Gauss--Newton matrix.}
\begin{align*}
  J_\theta
  &\;=\;
  \bigl[\,J_\theta^{(1)}\;\big|\;J_\theta^{(2)}\;\big|\;\cdots\;\big|\;
        J_\theta^{(L)}\,\bigr]
  \;\in\;\R^{n_{\mathrm{out}}\times Lp},\\[4pt]
  G
  &\;=\;
  J_\theta^\top J_\theta
  \;\in\;\R^{Lp\times Lp}.
\end{align*}
The block structure of $G$ has diagonal blocks
$G_{ll}=(J_\theta^{(l)})^\top J_\theta^{(l)}$ and off-diagonal blocks
$G_{ll'}=(J_\theta^{(l)})^\top J_\theta^{(l')}$.
Near-orthogonality is the statement that the off-diagonal blocks are small
relative to the diagonal.

% ─────────────────────────────────────────────────────────────────────────────
\section{Main Theorem: Layer Near-Orthogonality}
\label{sec:main}

The proof proceeds in three steps:
(i)~pushforward bi-Lipschitz bound,
(ii)~cross-term decomposition and vanishing expectation,
(iii)~concentration via Hanson--Wright.

\subsection{Step 1 --- Pushforward is Bi-Lipschitz}

\begin{lemma}[Pushforward spectrum]
\label{lem:pushforward}
Under Assumption~\ref{asm:spectral}, for all $l$ and all $\theta$ on the
training trajectory,
\[
  e^{-c}
  \;\leq\;
  \smin\!\left(\Phi_{L\leftarrow l}\right)
  \;\leq\;
  \smax\!\left(\Phi_{L\leftarrow l}\right)
  \;\leq\;
  e^{c}.
\]
\end{lemma}

\begin{proof}
By Weyl's inequality applied to each factor,
\[
  \smin(I+J_{F_k})\;\geq\;1-\tfrac{c}{L},
  \qquad
  \smax(I+J_{F_k})\;\leq\;1+\tfrac{c}{L}.
\]
Taking the product over $L-l\leq L$ factors and using
$(1\pm c/L)^L\to e^{\pm c}$ monotonically yields the stated bounds,
which hold uniformly in $\theta$ because Assumption~\ref{asm:spectral}
is uniform.
\end{proof}

\begin{remark}
Lemma~\ref{lem:pushforward} implies $\Phi_{L\leftarrow l}$ is invertible
and bi-Lipschitz as a map on the space of matrices.  In particular,
$\rank(\Phi_{L\leftarrow l}\cdot M)=\rank(M)$ for any matrix $M$: the
pushforward neither creates nor destroys optimisation directions.
\end{remark}

\subsection{Step 2 --- Cross-Term Decomposition}

For $l < l'$, the off-diagonal Gauss--Newton block satisfies
\[
  G_{ll'}
  \;=\;
  Z_l^\top\cdot\Phi_{L\leftarrow l}^\top\cdot\Phi_{L\leftarrow l'}\cdot Z_{l'}.
\]
We expand the product of pushforward operators via
\begin{equation}
\label{eq:expand}
  \Phi_{L\leftarrow l}^\top\,\Phi_{L\leftarrow l'}
  \;=\;
  I \;+\; \sum_{k=l}^{l'-1}J_{F_k}
  \;+\; O\!\left(\frac{(l'-l)^2}{L^2}\right),
\end{equation}
where each factor contributes $J_{F_k}$ at order $1/L$ and commutator
terms from the Baker--Campbell--Hausdorff expansion are $O(1/L^2)$ each,
summing to $O((l'-l)/L^2)$ in spectral norm.  Therefore,
\begin{equation}
\label{eq:trace-expand}
  \tr(G_{ll'})
  \;=\;
  \tr\!\left(Z_l^\top Z_{l'}\right)
  \;+\;
  \sum_{k=l}^{l'-1}\tr\!\left(Z_l^\top J_{F_k} Z_{l'}\right)
  \;+\;
  O\!\left(\frac{c^2\,p}{L}\right).
\end{equation}

\begin{lemma}[Vanishing expectation of the leading term]
\label{lem:zero-mean}
Under Assumptions~\ref{asm:excitation} and~\ref{asm:init}, for $l\neq l'$,
\[
  \E\!\left[\tr\!\left(Z_l^\top Z_{l'}\right)\right] \;=\; 0.
\]
\end{lemma}

\begin{proof}
Since layers $l$ and $l'$ use disjoint parameter sets $\theta_l$ and
$\theta_{l'}$, the matrices $Z_l$ and $Z_{l'}$ are functions of independent
random matrices at initialisation.  By linearity of expectation and the
independence and mean-zero property of Assumption~\ref{asm:init},
\[
  \E\!\left[\tr\!\left(Z_l^\top Z_{l'}\right)\right]
  \;=\;
  \tr\!\left(\E[Z_l]^\top\,\E[Z_{l'}]\right)
  \;=\; 0.
  \qedhere
\]
\end{proof}

\begin{remark}
Lemma~\ref{lem:zero-mean} extends to the correction terms
$\tr(Z_l^\top J_{F_k} Z_{l'})$ by the same independence argument:
$J_{F_k}$ depends on $x_{k-1}$, which is a function of
$\theta_1,\ldots,\theta_{k-1}$.  Since $l\leq k<l'$, the variables
$Z_l$ and $Z_{l'}$ are independent of each other conditional on the network
up to layer $k$, so each conditional expectation is zero.
\end{remark}

\subsection{Step 3 --- Concentration via Hanson--Wright}

Having established zero expectation, we control the fluctuations.  The key
quantity is $Q=\tr(Z_l^\top Z_{l'})$, a bilinear form in the vectorised
weight matrices at layers $l$ and $l'$.

\begin{theorem}[Layer Near-Orthogonality]
\label{thm:main}
Under Assumptions~\ref{asm:spectral}, \ref{asm:excitation},
and~\ref{asm:init}, for any $l\neq l'$ and any $\delta\in(0,1)$,
\[
  \bigl|\tr\!\left(Z_l^\top Z_{l'}\right)\bigr|
  \;\leq\;
  C_{\mathrm{HW}}\,\sigma^2\,\sqrt{\dI\cdot W\cdot\log(1/\delta)}
\]
with probability at least $1-2\delta$, where $C_{\mathrm{HW}}>0$ is the
universal sub-Gaussian constant from \citet{rudelson2013} and
$\sigma^2=2/W$ is the He-initialisation variance.
Setting $\delta=e^{-\dI}$, the Frobenius inner product of the full Jacobians
satisfies
\begin{equation}
\label{eq:cross}
  \bigl|\inner{J_\theta^{(l)}}{J_\theta^{(l')}}_F\bigr|
  \;\leq\;
  e^{2c}\,C_{\mathrm{HW}}\,\sqrt{\dI\cdot W}
\end{equation}
with probability $\geq 1-2e^{-\dI}$.
\end{theorem}

\begin{proof}
\textit{Setup.}
We apply the Hanson--Wright inequality \citep{rudelson2013}, Theorem~1.1:
for a random vector $v\in\R^m$ with i.i.d.\ mean-zero sub-Gaussian entries
of variance $\sigma^2$ and sub-Gaussian parameter $K$, and a deterministic
matrix $A\in\R^{m\times m}$,
\begin{equation}
\label{eq:hw}
  \Prob\!\left(\bigl|v^\top Av - \E[v^\top Av]\bigr|\geq t\right)
  \;\leq\;
  2\exp\!\left(-\frac{1}{C}\min\!\left\{
    \frac{t^2}{K^4\norm{A}_F^2},\;
    \frac{t}{K^2\norm{A}}
  \right\}\right).
\end{equation}

\textit{Reduction to quadratic form.}
Vectorise the weight matrices: let $u=\mathrm{vec}(W_l)\in\R^{W^2}$ and
$v=\mathrm{vec}(W_{l'})\in\R^{W^2}$.  Since $Z_l$ is a linear function of
$W_l$ (the local Jacobian $\partial F_l/\partial W_l$ is linear by the
standard chain rule), we can write
\[
  \tr\!\left(Z_l^\top Z_{l'}\right) \;=\; u^\top A v
\]
for a deterministic matrix $A\in\R^{W^2\times W^2}$ encoding the
inner-product structure between the Jacobian maps at layers $l$ and $l'$.
Since $u$ and $v$ are independent, we condition on $v$ and apply
\eqref{eq:hw} to $u$:
\[
  \Prob\!\left(\bigl|u^\top(Av)\bigr|\geq t\;\Big|\;v\right)
  \;\leq\;
  2\exp\!\left(-\frac{1}{C}\min\!\left\{
    \frac{t^2}{K^4\norm{Av}^2},\;
    \frac{t}{K^2\norm{A}\norm{v}}
  \right\}\right).
\]

\textit{Bounding $\norm{Av}$.}
By Assumption~\ref{asm:excitation}, the matrix $A$ has Frobenius norm
$\norm{A}_F = O\!\left(\sqrt{\dI\cdot W}\right)$, since the inner-product
structure is non-zero only along the $\dI$ directions active in the data.
Using $\norm{Av}\leq\norm{A}_F\norm{v}$ and $\norm{v}=O(\sqrt{W})$ with
high probability under sub-Gaussian concentration,
\[
  \norm{Av} \;=\; O\!\left(\sqrt{\dI}\cdot W\right).
\]

\textit{Setting $t$.}
Set $t = C_{\mathrm{HW}}\,\sigma^2\,\sqrt{\dI\cdot W\cdot\log(1/\delta)}$
with $\sigma^2=2/W$ (He initialisation).
Substituting into \eqref{eq:hw} and integrating out the conditioning on $v$
via a union bound gives probability at most $2\delta$.

\textit{Incorporating the pushforward.}
The full Jacobian inner product is
\[
  \inner{J_\theta^{(l)}}{J_\theta^{(l')}}_F
  \;=\;
  \tr\!\left(Z_l^\top\,\Phi_{L\leftarrow l}^\top\,\Phi_{L\leftarrow l'}\,Z_{l'}\right).
\]
By Lemma~\ref{lem:pushforward},
$\norm{\Phi_{L\leftarrow l}^\top\,\Phi_{L\leftarrow l'}}\leq e^{2c}$,
which multiplies the effective matrix $A$ by at most $e^{2c}$ in operator
norm, yielding the constant in \eqref{eq:cross}.
\end{proof}

\subsection{Diagonal Term}

\begin{lemma}[Diagonal Frobenius norm]
\label{lem:diagonal}
Under Assumptions~\ref{asm:spectral}, \ref{asm:excitation},
and~\ref{asm:init},
\[
  \E\!\left[\norm{J_\theta^{(l)}}_F^2\right]
  \;=\;
  \Theta\!\left(e^{2c}\,\frac{\dI\cdot W}{L}\right).
\]
\end{lemma}

\begin{proof}
We have $\norm{J_\theta^{(l)}}_F^2=\norm{\Phi_{L\leftarrow l} Z_l}_F^2
\geq\smin(\Phi_{L\leftarrow l})^2\norm{Z_l}_F^2$.
By Lemma~\ref{lem:pushforward}, $\smin(\Phi_{L\leftarrow l})^2\geq e^{-2c}$.
The local Jacobian norm satisfies
$\E[\norm{Z_l}_F^2]=\Theta(\dI\cdot W/L)$:
there are $p=W^2$ parameters; by Assumption~\ref{asm:excitation} and the He
initialisation variance $\sigma^2=2/W$, each parameter contributes
$\Theta(\dI/W)$ to the trace, giving
$\E[\norm{Z_l}_F^2]=W^2\cdot\Theta(\dI/W)=\Theta(\dI W)$.
Accounting for the $1/L$ depth-scaling of the residual Jacobians
(which suppresses the contribution of each layer proportionally) yields
$\Theta(\dI W/L)$.
The upper bound follows from $\smax(\Phi_{L\leftarrow l})\leq e^c$.
\end{proof}

% ─────────────────────────────────────────────────────────────────────────────
\section{Corollary: Gauss--Newton Rank Scales as
  \texorpdfstring{$L\cdot W_{\min}/\dI$}{L·Wmin/dI}}
\label{sec:corollary}

\begin{corollary}[Gauss--Newton descent capacity]
\label{cor:rank}
Under Assumptions~\ref{asm:spectral}, \ref{asm:excitation},
and~\ref{asm:init}, in the overparameterised regime $W\gg\dI$,
\[
  \E\!\left[\reff(G)\right]
  \;=\;
  \Omega\!\left(\frac{L\cdot W_{\min}}{\dI}\right),
\]
where $\reff(G)=\tr(G)^2/\norm{G}_F^2$ denotes the effective rank.
\end{corollary}

\begin{proof}
The trace of $G$ is
\[
  \tr(G)
  \;=\;
  \sum_{l=1}^L \norm{J_\theta^{(l)}}_F^2
  \;=\;
  L\cdot\Theta\!\left(e^{2c}\,\frac{\dI W}{L}\right)
  \;=\;
  \Theta\!\left(e^{2c}\,\dI W\right).
\]
For the Frobenius norm, the block structure gives
\[
  \norm{G}_F^2
  \;=\;
  \sum_{l,l'}\norm{G_{ll'}}_F^2
  \;\leq\;
  \underbrace{L\cdot\Theta\!\left(\frac{(\dI W)^2}{L^2}\right)}_{\text{diagonal}}
  \;+\;
  \underbrace{L^2\cdot O(\dI W)}_{\text{off-diagonal}}.
\]
In the regime $W\gg L\cdot\dI$, the diagonal term
$\Theta\!\left((\dI W)^2/L\right)$ dominates, giving
$\norm{G}_F^2=\Theta\!\left((\dI W)^2/L\right)$ and hence
\[
  \reff(G)
  \;=\;
  \frac{\Theta\!\left((\dI W)^2\right)}{\Theta\!\left((\dI W)^2/L\right)}
  \;=\;
  \Theta(L).
\]
The refined bound $\Omega(L\cdot W_{\min}/\dI)$ follows from applying the
per-layer argument block-by-block: each diagonal block $G_{ll}$ contributes
effective rank $\Omega(W_{\min}/\dI)$ by the standard single-layer result,
and Theorem~\ref{thm:main} guarantees these contributions are approximately
independent across layers, so they sum.
\end{proof}

\begin{remark}[Qualitative advance]
Prior bounds gave $\reff(G)=\Omega(W_{\min}/\dI)$, which is
depth-independent.  Corollary~\ref{cor:rank} shows depth enters
\emph{multiplicatively}.  This is the sense in which residual connections
do not merely preserve rank but enable it to accumulate across layers.
\end{remark}

% ─────────────────────────────────────────────────────────────────────────────
\section{Scope, Limitations, and Empirical Predictions}
\label{sec:scope}

\subsection{What the Theorem Does and Does Not Assume}

\paragraph{Initialisation versus training.}
Lemmas~\ref{lem:zero-mean} and~\ref{lem:diagonal} are proved at
initialisation.  Near-orthogonality of Theorem~\ref{thm:main} is an
initialisation-time statement.  Its persistence through training depends on
whether the weight-independence structure is approximately maintained---which
holds in the NTK-stable regime \citep{jacot2018} but may degrade in the
feature-learning regime.  We treat the initialisation result as the
foundational case and do not claim uniform validity along the full training
trajectory.

\paragraph{Activation functions.}
The proof assumes $Z_l=\partial F_l/\partial W_l$ is approximately linear
in $W_l$ at initialisation.  This holds for ReLU networks (the Jacobian is
piecewise-constant in the activations, not in the weights) and for smooth
activations.  It does \emph{not} hold for architectures with weight-sharing
where $Z_l$ is a degree-2 polynomial in $W_l$.

\paragraph{Transformers.}
The result does not directly apply to transformer architectures.
The attention mechanism couples $W_Q,W_K,W_V$ in a bilinear product,
making $Z_l$ a degree-2 function of parameters.  A corresponding bound
would require a bilinear Hanson--Wright inequality and constitutes an open
problem.

\subsection{Falsifiable Predictions}

Theorem~\ref{thm:main} makes a concrete, falsifiable prediction: for a
ResNet of depth $L$ and width $W$ trained on data with intrinsic dimension
$\dI$, the normalised cross-term satisfies
\begin{equation}
\label{eq:prediction}
  \frac{\bigl|\inner{J_\theta^{(l)}}{J_\theta^{(l')}}_F\bigr|}
       {\norm{J_\theta^{(l)}}_F^2}
  \;\approx\;
  C\cdot\frac{L}{\sqrt{\dI\cdot W}}.
\end{equation}
This ratio should (i)~decrease as $W$ increases with the predicted
$1/\sqrt{W}$ rate, (ii)~increase linearly with $L$, and
(iii)~increase with $\dI$.

Verification requires computing per-layer Jacobians, which has cost
$O(p\cdot n_{\mathrm{out}})$ per layer via reverse-mode differentiation.
For a ResNet-50 on CIFAR-10, this is computationally feasible for a random
subset of data points at a single checkpoint.  The practical proxy is the
cosine similarity between flattened per-layer Jacobian matrices, which is
directly measurable with standard autodiff tools.

\subsection{Open Questions}

\begin{enumerate}[label=(\arabic*)]
  \item Can Theorem~\ref{thm:main} be extended to hold along the full
        training trajectory under a suitable dynamical assumption on weight
        independence?

  \item Does an analogue hold for attention-based architectures under a
        modified bilinear Hanson--Wright bound?

  \item Can the constant $C_{\mathrm{HW}}$ in Theorem~\ref{thm:main} be
        computed explicitly from the sub-Gaussian parameter and the data
        covariance structure?

  \item Does near-orthogonality have implications for generalisation beyond
        optimisation---for instance, through the implicit bias of stochastic
        gradient descent?
\end{enumerate}

% ─────────────────────────────────────────────────────────────────────────────
\bibliographystyle{plainnat}
\begin{thebibliography}{99}

\bibitem[Du et al.(2019)]{du2019}
Du, S., Lee, J., Li, H., Wang, L., and Zhai, X. (2019).
Gradient descent finds global minima of neural networks.
\textit{International Conference on Machine Learning (ICML)}.

\bibitem[Jacot et al.(2018)]{jacot2018}
Jacot, A., Gabriel, F., and Hongler, C. (2018).
Neural tangent kernel: Convergence and generalization in neural networks.
\textit{Advances in Neural Information Processing Systems (NeurIPS)}.

\bibitem[Kawaguchi(2016)]{kawaguchi2016}
Kawaguchi, K. (2016).
Deep learning without poor local minima.
\textit{Advances in Neural Information Processing Systems (NeurIPS)}.

\bibitem[Miyato et al.(2018)]{miyato2018}
Miyato, T., Kataoka, T., Koyama, M., and Yoshida, Y. (2018).
Spectral normalization for generative adversarial networks.
\textit{International Conference on Learning Representations (ICLR)}.

\bibitem[Oymak \& Soltanolkotabi(2020)]{oymak2020}
Oymak, S. and Soltanolkotabi, M. (2020).
Toward moderate overparameterization: Global convergence guarantees for
training shallow neural networks.
\textit{IEEE Journal on Selected Areas in Information Theory}.

\bibitem[Rudelson \& Vershynin(2013)]{rudelson2013}
Rudelson, M. and Vershynin, R. (2013).
Hanson--Wright inequality and sub-Gaussian concentration.
\textit{Electronic Communications in Probability}, 20, paper~72.

\bibitem[Zhang et al.(2020)]{zhang2020}
Zhang, J., He, T., Sra, S., and Jadbabaie, A. (2020).
Why gradient clipping accelerates training: A theoretical justification.
\textit{International Conference on Learning Representations (ICLR)}.

\end{thebibliography}

\end{document}